\section{Алгоритмическая реализация Time-Dependent A* и SIMD-оптимизации}

Сердцем модуля \texttt{router} является класс \texttt{TdAltRouter}, реализующий алгоритм поиска кратчайшего пути с временной зависимостью (Time-Dependent A*). Для обеспечения заявленной пропускной способности в тысячи запросов в секунду код маршрутизатора глубоко оптимизирован с использованием векторных инструкций процессора.

\subsection{Векторизация 4-арной кучи (AVX2)}

В классической реализации алгоритма Дейкстры или A* основным узким местом (занимающим до 60\% процессорного времени) является извлечение узла с минимальным весом из очереди с приоритетом (Priority Queue). Традиционно для этого используется бинарная куча (\texttt{std::priority\_queue}), имеющая глубину $\log_2(V)$.

В разработанном ядре осуществлен переход на \textbf{4-арную кучу} (4-ary Heap), которая уменьшает глубину дерева в два раза ($\log_4(V)$), что критически важно для снижения количества промахов кэша. Однако в 4-арной куче операция проталкивания узла вниз (Sift-down) требует нахождения минимума среди 4 потомков на каждом уровне, что увеличивает количество инструкций сравнения и ветвлений.

Для нивелирования этих затрат в проекте реализован поиск минимального потомка с использованием 256-битных векторных инструкций AVX2 (Advanced Vector Extensions).
Четыре 64-битные записи потомков (содержащие вес и ID узла) загружаются в один YMM-регистр:
\begin{verbatim}
inline int find_min_child_avx2(const uint32_t* data_ptr) {
    __m256i data = _mm256_loadu_si256(reinterpret_cast<const __m256i*>(data_ptr));
    // Наложение маски для извлечения 32-битных весов из 64-битных записей
    __m256i id_mask_and = _mm256_set1_epi64x(0xFFFFFFFF00000000);
    __m256i weights_only = _mm256_and_si256(data, id_mask_and);
    
    // Поиск минимального значения (через перестановки) и генерация маски
    // ...
    int child_idx = __builtin_ctz(bitmask) / 8;
    return child_idx;
}
\end{verbatim}
Векторизация устраняет все условные переходы (\texttt{if-else}) при поиске дочернего узла. Процессор выбирает минимального потомка за несколько аппаратных тактов без риска сброса конвейера из-за ошибочного предсказания переходов (Branch Misprediction).

\subsection{Механизм Lock-Free резервирования и интерполяции (LERP)}

В процессе раскрытия узлов (Expand) алгоритм TD-A* оценивает вес каждого маневра. Для чтения объемов трафика из \texttt{volume\_buckets} маршрутизатор не использует транзакционные мьютексы. Массивы атомарных переменных считываются со слабым порядком памяти (\texttt{memory\_order\_relaxed}). 

Для устранения ступенчатости штрафов реализован механизм линейной интерполяции (LERP) времени, который считывает объемы из двух соседних 5-минутных слотов ($P_1$ и $P_2$) и интерполирует непосредственно квадраты плотности потока, применяя заранее вычисленный \texttt{K\_magic}:
\begin{verbatim}
uint32_t p1 = K_magic * v1 * v1;
uint32_t p2 = K_magic * v2 * v2;
uint32_t bpr_penalty = (p1 + (p2 - p1) * (eta % 300) / 300) >> 20;
\end{verbatim}
После того как оптимальный путь (цепочка сегментов) найден, роутер осуществляет резервирование емкости. Проходя по восстановленному пути, для каждого сегмента вызывается операция \texttt{fetch\_add}, которая атомарно инкрементирует значение в соответствующей временной корзинке. Это математически аппроксимирует эффект перераспределения потока (Congestion Pricing), уводя следующих агентов на альтернативные маршруты.

\subsection{Пространственные индексы МПР (In-Memory SpatialGrid)}

Для обеспечения мгновенной реакции Модуля принятия решений (МПР) на инциденты реализован сверхбыстрый пространственный индекс. Классические деревья (R-Tree), основанные на указателях в куче, были заменены на плоскую регулярную сетку (SpatialGrid), развернутую в непрерывном массиве. 

Привязка географической геометрии дорог к ячейкам сетки выполняется на этапе офлайн-препроцессинга. В реальном времени МПР формирует ограничивающий прямоугольник (Bounding Box) зоны ДТП или перекрытия и за $O(1)$ получает доступ к списку ID ребер, попавших в ячейки. Затем МПР выставляет директивные штрафы в атомарное поле \texttt{mpr\_penalty} целевых сегментов. Благодаря строгой архитектуре кэширования, поиск затронутых ребер и обновление штрафов происходят за доли микросекунды, не блокируя работу соседних потоков пула маршрутизатора.
